\documentclass[11pt]{article}

\usepackage[margin=1in]{geometry}
\usepackage{amsmath,amssymb,amsthm}
\usepackage{booktabs}
\usepackage{hyperref}

\newtheorem{theorem}{Theorem}
\newtheorem{lemma}[theorem]{Lemma}
\newtheorem{proposition}[theorem]{Proposition}
\newtheorem{corollary}[theorem]{Corollary}
\theoremstyle{remark}
\newtheorem{remark}[theorem]{Remark}

\title{Normalising the Force by $|\nabla\rho|$:\\
       Why it Removes the Dip}
\author{}
\date{}

\begin{document}
\maketitle

\begin{abstract}
An elongated (elliptical) potential field produces a spurious local minimum
of $|F|$ along its major axis --- weakest force exactly where the obstacle
lies most directly ahead. This note explains, first informally and then with
proofs, why dividing the force by $|\nabla\rho|$ removes that minimum for
every elongation while leaving the force direction exactly unchanged. The
essential observation is that $|\nabla U|$ factors into a part that depends
on $\rho$ alone and a part that depends on \emph{direction}, and that the
entire defect lives in the second part. We also delimit the result: it is
exact on horizontal scans of the implemented field, but the stronger
statement that $|F|$ is constant on ellipses fails there, for reasons we make
precise.
\end{abstract}

\section{The idea in words}
\label{sec:intuition}

\subsection{Force is contour spacing}

The magnitude of a gradient is a rate: how much $U$ changes per metre
travelled. Equivalently, it measures how tightly the equipotential contours
are packed. Where contours crowd, the force is large; where they spread, it
is small.

Now elongate the field by a factor $q$ along $y$. The contours become
ellipses stretched $q:1$. Crucially, \emph{the same potential drop is now
spread over $q$ times more metres at the poles of the ellipse than at its
flanks}. So the force at the poles is $q$ times weaker --- not because the
situation there is any safer, but because of how the contours were drawn. The
dip is a measurement artefact of this stretching.

\subsection{$\rho$ is a risk coordinate, not a length}

The cleanest way to see the artefact is to notice that $\rho$ is not a
distance. It is a \emph{risk coordinate}: a bookkeeping variable in which
one unit longitudinally is worth $t_y$ metres while one unit laterally is
worth only $t_x$ metres. When we form $\nabla U$ we are implicitly applying
the chain rule,
\begin{equation}
  \underbrace{\frac{\partial U}{\partial(\text{metres})}}_{\text{what }|\nabla U|\text{ measures}}
  \;=\;
  \underbrace{\frac{\partial U}{\partial\rho}}_{\text{what we want}}
  \;\times\;
  \underbrace{\frac{\partial\rho}{\partial(\text{metres})}}_{\text{exchange rate}} ,
\end{equation}
and the trouble is that the exchange rate is \emph{direction-dependent}: it
is $1/t_y$ along the major axis and $1/t_x$ along the minor axis. The first
factor is the honest quantity --- how fast risk changes per unit of risk ---
and it can have no directional preference, because $U$ is a function of
$\rho$ alone. The second factor is pure bookkeeping, and it is the sole
source of the dip.

Dividing by $|\nabla\rho|$ cancels the bookkeeping factor and returns
$\partial U/\partial\rho$. That is the whole method.

\subsection{The eikonal picture}

There is a standard name for the property $\rho$ lacks. A genuine distance
function satisfies the \emph{eikonal condition}
\begin{equation}
  \bigl|\nabla\phi\bigr| \;=\; 1 ,
\end{equation}
meaning one metre of travel buys exactly one unit of $\phi$, in every
direction. Our $\rho$ violates this: its gradient norm ranges over
$[1/t_y,\,1/t_x]$, a swing of exactly the elongation ratio $q$. In level-set
methods a function that has drifted away from the eikonal condition is
repaired by \emph{reinitialisation}; dividing by $|\nabla\rho|$ is the
closed-form, one-step version of the same repair. Seen this way, the method
is not a trick but the restoration of a property the field was implicitly
assumed to have.

\subsection{Why the direction survives}

Finally, the reason this costs nothing in steering behaviour: $|\nabla\rho|$
is a \emph{positive scalar}. Dividing a vector by a positive scalar changes
its length and not its heading. So every arrow in the field keeps pointing
exactly where it pointed before; only the colour map changes. This is what
distinguishes the method from metric-weighted alternatives such as
$-G^{-1}\nabla U$, which rotate the force toward the radial direction and
destroy the lateral guidance.

\begin{remark}[One-sentence summary]
The dip is the gap between ``potential drop per metre'' and ``potential drop
per unit of risk''; the normalisation measures the latter, which by
construction cannot prefer a direction.
\end{remark}

\section{Setup and notation}

Let the obstacle sit at $(x_{o},y_{o})$ and write the gaps
$d_{x}=x_{o}-X$, $d_{y}=y_{o}-Y$ for a host at $(X,Y)$. With per-axis time
constants $t_{x},t_{y}>0$ define the elliptical pseudo-distance
\begin{equation}
  \rho \;=\; \sqrt{\Bigl(\frac{d_{x}}{t_{x}}\Bigr)^{2}
                 + \Bigl(\frac{d_{y}}{t_{y}}\Bigr)^{2}},
  \qquad
  q \;=\; \frac{t_{y}}{t_{x}}\ \ (\text{the elongation}).
  \label{eq:rho}
\end{equation}
Throughout, $\nabla$ is taken with respect to the host position $(X,Y)$;
since $d_{x},d_{y}$ are affine in $X,Y$ with unit slopes, signs cancel in
every magnitude below. We consider potentials of the form $U=U(\rho)$ with
$U'$ of one sign, and define the normalised force
\begin{equation}
  \boxed{\;F_{\mathrm{n}} \;=\; -\,\frac{\nabla U}{\bigl|\nabla\rho\bigr|}\;}
  \label{eq:Fn}
\end{equation}

\section{The mathematics}

\subsection{Localising the defect}

\begin{lemma}[Factorisation]
\label{lem:factor}
If $U=U(\rho)$ then $\nabla U = U'(\rho)\,\nabla\rho$, and hence
\begin{equation}
  \bigl|\nabla U\bigr|
  \;=\;
  \underbrace{\bigl|U'(\rho)\bigr|}_{\text{depends on }\rho\text{ only}}
  \cdot
  \underbrace{\bigl|\nabla\rho\bigr|}_{\text{depends on direction}} .
  \label{eq:factor}
\end{equation}
\end{lemma}

\begin{proof}
Immediate from the chain rule applied to the composition $U\circ\rho$.
\end{proof}

The two factors have sharply different characters, and the next two results
make that precise. Lemma~\ref{lem:factor} is the formal content of
Section~\ref{sec:intuition}: the first factor is the honest rate, the second
is the exchange rate.

\begin{lemma}[Range of the exchange rate]
\label{lem:range}
Away from the obstacle centre,
\begin{equation}
  \bigl|\nabla\rho\bigr|
  =\frac{\sqrt{\bigl(d_{x}/t_{x}^{2}\bigr)^{2}+\bigl(d_{y}/t_{y}^{2}\bigr)^{2}}}{\rho},
  \qquad
  \bigl|\nabla\rho\bigr|\in
  \Bigl[\tfrac{1}{\max(t_{x},t_{y})},\ \tfrac{1}{\min(t_{x},t_{y})}\Bigr],
\end{equation}
with the endpoints attained exactly on the two axes: $|\nabla\rho|=1/t_{y}$
on the major axis $(d_{x}=0)$ and $|\nabla\rho|=1/t_{x}$ on the minor axis
$(d_{y}=0)$. The ratio between the extremes is exactly $q$.
\end{lemma}

\begin{proof}
Differentiating \eqref{eq:rho} gives
$\partial\rho/\partial d_{x}=d_{x}/(t_{x}^{2}\rho)$ and likewise for
$d_{y}$, which yields the stated formula. For the range, put
$u=(d_{x}/t_{x})^{2}\ge0$ and $v=(d_{y}/t_{y})^{2}\ge0$, so that
$\rho^{2}=u+v$ and
$(d_{x}/t_{x}^{2})^{2}+(d_{y}/t_{y}^{2})^{2}=u/t_{x}^{2}+v/t_{y}^{2}$.
Hence
\begin{equation}
  \bigl|\nabla\rho\bigr|^{2}
  =\frac{u}{u+v}\cdot\frac{1}{t_{x}^{2}}
  +\frac{v}{u+v}\cdot\frac{1}{t_{y}^{2}},
\end{equation}
a convex combination of $1/t_{x}^{2}$ and $1/t_{y}^{2}$ with non-negative
weights summing to one. A convex combination lies between its extremes, and
attains them precisely when a weight vanishes, i.e.\ when $v=0$ or $u=0$.
Taking square roots gives the claim.
\end{proof}

\noindent
Lemma~\ref{lem:range} quantifies the intuition exactly: the offending factor
sweeps a range whose width \emph{is} the elongation $q$. For an isotropic
field $t_{x}=t_{y}$ the range collapses to a point and there is nothing to
correct.

\subsection{What the normalisation achieves}

\begin{theorem}[The normalised magnitude is direction-free]
\label{thm:magnitude}
For $U=U(\rho)$ and $F_{\mathrm{n}}$ as in \eqref{eq:Fn},
\begin{equation}
  \bigl|F_{\mathrm{n}}\bigr| \;=\; \bigl|U'(\rho)\bigr| ,
\end{equation}
a function of $\rho$ alone. Equivalently, $F_{\mathrm{n}}=-U'(\rho)\,\hat n$
with $\hat n=\nabla\rho/|\nabla\rho|$ the outward unit normal to the ellipse
through the point.
\end{theorem}

\begin{proof}
By Lemma~\ref{lem:factor}, $|\nabla U|=|U'(\rho)|\,|\nabla\rho|$. Dividing by
$|\nabla\rho|>0$ gives $|F_{\mathrm{n}}|=|U'(\rho)|$. The vector form follows
from $\nabla U=U'(\rho)\nabla\rho$ on dividing by $|\nabla\rho|$.
\end{proof}

\begin{corollary}[The two families of level sets coincide]
\label{cor:levelsets}
Under the hypotheses of Theorem~\ref{thm:magnitude}, the level sets of
$|F_{\mathrm{n}}|$ are exactly the level sets of $U$, namely the ellipses
$\rho=\text{const}$.
\end{corollary}

\begin{proof}
Both $|F_{\mathrm{n}}|=|U'(\rho)|$ and $U=U(\rho)$ are functions of $\rho$
alone, so each is constant on $\{\rho=c\}$; as $|U'|$ is of one sign and
$\rho$ ranges over an interval, the sets $\{\rho=c\}$ are precisely the level
sets of each.
\end{proof}

Corollary~\ref{cor:levelsets} is the sharpest statement of why the method
works. The dip was, by definition, a mismatch between the shape of the
$|F|$ contours (pinched, non-convex, peanut-like) and the shape of the $U$
contours (nested convex ellipses). The normalisation forces the two families
to agree, and a family of nested ellipses has no pinch.

\begin{theorem}[No dip, at any elongation]
\label{thm:nodip}
Fix a height $s>0$ and suppose $|U'|$ is strictly decreasing in $\rho$
(as for $U=\lambda/\rho$). Then along the horizontal scan $y=s$ the magnitude
$|F_{\mathrm{n}}|$ is strictly decreasing in $|x|$, attains its maximum at
$x=0$, and has no interior critical point. This holds for every $q>0$, with
no threshold.
\end{theorem}

\begin{proof}
By Theorem~\ref{thm:magnitude}, $|F_{\mathrm{n}}|=|U'(\rho)|$. At fixed $s$,
\begin{equation}
  \rho(x,s)=\sqrt{\Bigl(\frac{x}{t_{x}}\Bigr)^{2}+\Bigl(\frac{s}{t_{y}}\Bigr)^{2}}
\end{equation}
is strictly increasing in $|x|$, since $\partial\rho/\partial x = x/(t_{x}^{2}\rho)>0$
for $x>0$. The composition of the strictly decreasing map $\rho\mapsto|U'(\rho)|$
with the strictly increasing map $|x|\mapsto\rho$ is strictly decreasing.
A strictly monotone function on $[0,\infty)$ attains its maximum at the left
endpoint and has no interior stationary point. No hypothesis on $q$ was used.
\end{proof}

\begin{corollary}[Direction is exactly preserved]
\label{cor:direction}
$F_{\mathrm{n}}$ is everywhere parallel to $-\nabla U$, with the same
orientation.
\end{corollary}

\begin{proof}
$F_{\mathrm{n}} = (-\nabla U)\big/|\nabla\rho|$ and $|\nabla\rho|>0$ away from
the obstacle centre. Multiplication by a positive scalar preserves both the
line spanned by a vector and its orientation.
\end{proof}

\begin{proposition}[Consistency: the isotropic case is a fixed point]
\label{prop:isotropic}
If $t_{x}=t_{y}=t$ then $|\nabla\rho|\equiv 1/t$, so $F_{\mathrm{n}}$ differs
from $-\nabla U$ by the single global constant $t$; if moreover $t=1$ then
$F_{\mathrm{n}}=-\nabla U$ identically.
\end{proposition}

\begin{proof}
Set $t_{x}=t_{y}=t$ in Lemma~\ref{lem:range}: the interval degenerates to the
point $1/t$.
\end{proof}

\noindent
Proposition~\ref{prop:isotropic} is the sanity check that the method is a
genuine generalisation: it does nothing precisely in the case that had
nothing wrong with it, and its strength scales with the elongation that
caused the problem.

\section{Scope: the implemented field}
\label{sec:scope}

The implemented potential is \emph{not} of the form $U(\rho)$. It carries a
vehicle-dimension prefactor,
\begin{equation}
  U_{\mathrm{impl}} \;=\; \frac{\lambda}{\bigl|k'\bigr|},
  \qquad
  k' \;=\; c(d_{x},d_{y})\,\rho,
  \qquad
  c=\begin{cases}
      1-w/|d_{x}| & \text{region 1/3},\\
      1-l/|d_{y}| & \text{region 2/4},
    \end{cases}
  \label{eq:impl}
\end{equation}
and $c$ is not a function of $\rho$. Two consequences, and it is worth being
precise about which results survive.

\begin{proposition}[Theorem~\ref{thm:nodip} is exact on horizontal scans]
\label{prop:horizontal}
Restrict \eqref{eq:impl} to a horizontal line $Y=\text{const}$ lying in
region 2/4. Then $d_{y}$ is constant along that line, hence so is
$c=1-l/|d_{y}|$, and $U_{\mathrm{impl}}=\bigl(\lambda/c\bigr)\rho^{-1}$ is a
function of $\rho$ alone \emph{along that line}. Theorem~\ref{thm:nodip}
therefore applies verbatim, and the contrast satisfies $C=1$ exactly.
\end{proposition}

\begin{proof}
Along $Y=\text{const}$ we have $d_{y}=y_{o}-Y$ fixed, so $c$ is a constant
$c_{0}$ and $U_{\mathrm{impl}}=(\lambda/c_{0})\rho^{-1}$, to which
Theorem~\ref{thm:nodip} applies with $U'(\rho)=-(\lambda/c_{0})\rho^{-2}$,
strictly decreasing in magnitude.
\end{proof}

\begin{remark}[Corollary~\ref{cor:levelsets} does \emph{not} survive]
\label{rem:notlevelsets}
Along an ellipse $\rho=\text{const}$ the gap $d_{y}$ varies, so $c$ varies and
$|F_{\mathrm{n}}|$ is not constant there. Measured on the implemented field
at $q=15$, the coefficient of variation of $|F_{\mathrm{n}}|$ along arcs of
fixed $\rho$ is $52\%$ to $104\%$ --- reduced from $105\%$ to $162\%$ for the
unnormalised force, but far from zero. The guarantee delivered in practice is
thus Proposition~\ref{prop:horizontal} (no dip on any horizontal scan), not
the stronger Corollary~\ref{cor:levelsets}.
\end{remark}

\section{Numerical verification}

All figures below are measured on the implemented field at the current
configuration $t_{x}=1$, $t_{y}=15$, i.e.\ $q=15$.

\begin{center}
\begin{tabular}{rccc}
\toprule
$Y$ & plain $C$ & ridge $x^{*}$ & normalised $C$\\
\midrule
$36.5$ & $1.926$ & $0.288$ & $\mathbf{1.000}$\\
$37.5$ & $2.418$ & $0.340$ & $\mathbf{1.000}$\\
$38.5$ & $2.804$ & $0.390$ & $\mathbf{1.000}$\\
$39.5$ & $3.113$ & $0.438$ & $\mathbf{1.000}$\\
\bottomrule
\end{tabular}
\end{center}

\noindent
Lemma~\ref{lem:range} predicts $|\nabla\rho|\in[1/t_{y},1/t_{x}]=[0.0667,1]$;
sampled over the plotting window the measured range is $[0.0712,\,1.0000]$,
the lower end approached as the grid nears the major axis. The analytic
expression of Lemma~\ref{lem:range} agrees with a central-difference
evaluation to $6.4\times10^{-10}$. Corollary~\ref{cor:direction} is confirmed
to machine precision: the force angle at sample points above and below the
obstacle is unchanged to $0.00\times10^{0}$ degrees.

\section{The cost}

$F_{\mathrm{n}}$ is not a gradient field. Writing
$F_{\mathrm{n}}=-\Phi\,\nabla\rho$ with $\Phi=U'(\rho)/|\nabla\rho|$, such a
field is conservative only if $\Phi$ is a function of $\rho$ alone; but by
Lemma~\ref{lem:range} the factor $|\nabla\rho|$ varies along a single ellipse
$\rho=\text{const}$ whenever $q>1$, so $\Phi$ is not, and
$\nabla\times F_{\mathrm{n}}\neq0$. Consequently $U$ is no longer a Lyapunov
function for the normalised force. The magnitude of the violation, and the
comparison against alternative remedies, are treated in the companion note
\emph{Remedies for the Force-Magnitude Dip in an Inverse-Elliptical Potential
Field}.

\section{Relation to existing work}

The manoeuvre has two well-established relatives. In level-set methods, a
function that has lost the eikonal property $|\nabla\phi|=1$ is restored by
\emph{reinitialisation}, classically by evolving
$\partial_{\tau}\phi+\operatorname{sgn}(\phi_{0})(|\nabla\phi|-1)=0$;
\eqref{eq:Fn} is the closed-form one-step analogue. In potential-field
control, the fully normalised gradient $-\nabla\gamma/\lVert\nabla\gamma\rVert$
is used to decouple the commanded direction from a separately regulated
magnitude, and incurs the same loss of conservativeness. Equation
\eqref{eq:Fn} sits between the two: it removes only the metric factor
$|\nabla\rho|$, retaining the physically meaningful variation $|U'(\rho)|$
that a full normalisation would discard.

\end{document}
